\chapter{Problem Analysis}
\label{chap:analysis}

The analysis of the problem to be solved is the basis of every engineering
like procedure. In this chapter, the overall problem is decomposed into sub
problems on the basis of the fundamentals of Chapter~2, the existing
solution approaches from the literature are analysed, and the requirements
for the system to be developed are derived.

\section{Decomposition of the Problem}
\label{sec:decomposition}

The overall problem, learning a sequence of tasks without losing old
knowledge, can be decomposed into the following sub problems:

\begin{enumerate}[label=\textbf{SP\arabic*.}, leftmargin=3em]
  \item \textbf{Representation of knowledge.} In which part of the network
        is the knowledge about a task stored, and how can new knowledge be
        added without changing the old one?
  \item \textbf{Assignment of tasks to parameters.} Which parameters should
        be updated when a new task arrives, and who makes this decision?
  \item \textbf{Measurement of forgetting.} How can forgetting be measured
        objectively, so that different solutions can be compared fairly?
  \item \textbf{Fair comparison.} How can it be guaranteed that an observed
        improvement comes from the investigated mechanism and not, for
        example, simply from a larger number of parameters?
\end{enumerate}

These sub problems depend on each other. The assignment of tasks to
parameters (SP2) decides directly about the representation of knowledge
(SP1): if two tasks share all parameters, interference is unavoidable, as
explained in Section~2.3. The measurement (SP3) is needed before any
solution can be judged, and the fairness question (SP4) decides which
baseline models the experiment needs. The design in Chapter~4 answers SP1
and SP2 with the Mixture of Experts architecture and its router, while the
experiment protocol in Chapters~6 and~7 answers SP3 and SP4.

\section{Analysis of Existing Solution Approaches}
\label{sec:existing}

Continual learning has been studied intensively in the last years, and
several surveys give a systematic overview
\cite{delange2021,wang2024survey,masana2022}. The existing methods can be
divided into three families.

\subsection{Regularisation Based Methods}

Regularisation based methods leave the network architecture unchanged and
modify the loss function instead. The most influential representative is
Elastic Weight Consolidation (EWC) \cite{kirkpatrick2017}. After finishing
a task, EWC estimates the importance of every parameter by the diagonal of
the Fisher information matrix,
\begin{equation}
  F_{ii} = \mathbb{E}\!\left[\left(
    \frac{\partial \log p(y \mid x; \theta)}{\partial \theta_i}
  \right)^{2}\right],
  \label{eq:fisher}
\end{equation}
and afterwards punishes quadratic deviations of important parameters from
their stored values $\theta_i^{*}$:
\begin{equation}
  \mathcal{L}_{\mathrm{EWC}}(\theta)
    = \mathcal{L}_{t}(\theta)
      + \frac{\lambda}{2} \sum_{i} F_{ii}\,(\theta_i - \theta_i^{*})^{2}.
  \label{eq:ewc}
\end{equation}
The coefficient $\lambda$ controls the stability plasticity trade off.
Related methods estimate the importance differently: Synaptic Intelligence
accumulates an online importance measure along the optimisation path
\cite{zenke2017}, and Memory Aware Synapses estimates the importance from
the sensitivity of the output function \cite{aljundi2018}. A shared
limitation is that the quadratic penalty is only a local approximation:
with more and more tasks, the collected constraints can rigidify the
network and finally prevent it from learning new tasks at all. A second
branch of regularisation avoids storing old data by distillation, for
example Learning without Forgetting \cite{li2017}.

\subsection{Replay Based Methods}

Replay methods keep a small episodic memory of real or generated samples
from earlier tasks and mix them into the training of new tasks. iCaRL
\cite{rebuffi2017} combines an exemplar memory with a nearest mean
classifier. Gradient Episodic Memory \cite{lopezpaz2017} uses the memory to
constrain the gradient updates so that the loss on every remembered task
cannot increase. Generative replay replaces the stored samples with data
from a learned generative model \cite{shin2017}. Replay belongs to the
strongest strategies, but it needs to store data, which can be impossible
for privacy reasons, and the memory grows with the number of tasks.

\subsection{Architecture Based Methods}

The third family changes the network itself. Progressive Neural Networks
\cite{rusu2016} allocate a new network column per task and connect it to
the frozen columns of earlier tasks; forgetting is impossible by
construction, but the number of parameters grows linearly with the number
of tasks, and the task identity must be known at test time. Other
approaches assign a binary mask to every task and train only the unmasked
weights \cite{mallya2018}. Expert Gate \cite{aljundi2017} is conceptually
closest to this project: it also uses a network of experts with a gating
mechanism, but it trains one new expert per task and selects experts by the
task identity, while the approach here lets the router learn the assignment
fully automatically without any task labels.

\subsection{Comparison of the Method Families}

Table~3.1 contrasts the three families along the criteria that matter for
this project. No family dominates the others. Regularisation methods are
cheap and need neither stored data nor task labels, but their protection is
only a local approximation of the true loss landscape, and the accumulated
constraints of many tasks can finally block further learning. Replay
methods reach the strongest results in most benchmarks, but they store old
data or a generative model of old data, which can be impossible under
privacy constraints, and their memory requirement grows with the length of
the task sequence. Architectural methods avoid both problems, but the
existing representatives either grow the network with every task or need
the task identity at test time to select the right component.

\begin{table}[H]
  \centering
  \small
  \caption{Comparison of the three continual learning method families.}
  \label{tab:families}
  \begin{tabular}{lccc}
    \toprule
    \textbf{Criterion} & \textbf{Regularisation} & \textbf{Replay} & \textbf{Architecture} \\
    \midrule
    Stores old data          & no   & yes  & no \\
    Needs task identity      & no   & no   & often \\
    Capacity growth per task & none & memory only & often linear \\
    Protection mechanism     & approximate & strong & structural \\
    Representative methods   & EWC, SI, MAS & iCaRL, GEM & PNN, Expert Gate \\
    \bottomrule
  \end{tabular}
\end{table}

For the present project, the decisive criterion is the combination of the
first three rows: a method that stores no data, needs no task labels and
does not grow without limit. Only a softly partitioned architecture can
offer all three properties at the same time. A Mixture of Experts model
with a learned router is a natural candidate, because its capacity is
fixed in advance but is used conditionally, so new tasks can occupy experts
that earlier tasks rarely touch.

\subsection{Conclusion for this Project}

The analysis shows a gap: regularisation methods are efficient but only
approximate, replay methods are strong but need stored data, and existing
architectural methods either grow without limit or need the task identity.
A Mixture of Experts model with a learned router promises a middle way: no
stored data, no task labels, and a capacity that is partitioned softly
instead of being fixed by hand. This promise is tested in the present
project.

\section{Requirements Analysis}
\label{sec:requirements}

From the problem analysis, the following requirements for the system were
collected and prioritised in four categories, as it is common practice for
project work.

\paragraph{Must have requirements.}
\begin{itemize}[leftmargin=2em]
  \item \textbf{M1:} Sequential training of a model on a sequence of tasks
        without access to the data of earlier tasks.
  \item \textbf{M2:} At least two continual learning benchmarks built from
        MNIST (SplitMNIST and PermutedMNIST).
  \item \textbf{M3:} A working Mixture of Experts model with a trainable
        router and more than two experts.
  \item \textbf{M4:} A naive baseline with the same training protocol for
        comparison.
  \item \textbf{M5:} Quantitative evaluation with the standard metrics
        average accuracy, average forgetting and backward transfer after
        every training stage.
\end{itemize}

\paragraph{Should have requirements.}
\begin{itemize}[leftmargin=2em]
  \item \textbf{S1:} A regularisation baseline (EWC) to position the MoE
        results relative to an established method.
  \item \textbf{S2:} A capacity matched wide baseline to separate capacity
        effects from routing effects (answers SP4).
  \item \textbf{S3:} Reproducibility through fixed random seeds and
        repeated runs.
  \item \textbf{S4:} Analysis of the router (expert usage per task).
\end{itemize}

\paragraph{Could have requirements.}
\begin{itemize}[leftmargin=2em]
  \item \textbf{C1:} An auxiliary loss for load balancing that keeps all
        experts active.
  \item \textbf{C2:} Configurable hyperparameters (number of experts, top
        $k$, epochs, learning rate).
\end{itemize}

\paragraph{Nice to have requirements.}
\begin{itemize}[leftmargin=2em]
  \item \textbf{N1:} Extension to further benchmarks or to the class
        incremental scenario.
  \item \textbf{N2:} Dynamic addition of new experts per task.
\end{itemize}

All must have and should have requirements were realised, and C1 and C2
were realised as well. N1 and N2 remain future work and are discussed in
Section~8.2.
